% !TeX root = ../navier-stokes-manuscript.tex
% ============================================================
% 2.1 Vortex Stretching
% ============================================================

The momentum equation has four pieces:
\[
  \underbrace{\partial_t u}_{\text{local acceleration}}
  +
  \underbrace{(u\cdot\nabla)u}_{\text{nonlinear transport}}
  =
  \underbrace{-\nabla p}_{\text{pressure}}
  +
  \underbrace{\nu\Delta u}_{\text{viscous diffusion}}.
\]
The central regularity problem lies in the competition between nonlinear transport and viscosity.
The hope is simple: viscosity smooths the fluid faster than nonlinear transport can sharpen it.
If that remains true at every scale, the velocity stays smooth.

\subsection{Vortex Stretching}\label{sub:ThePerfectlyStretchingVortex}

Picture one material vortex tube: a narrow column of rotating fluid marked by the particles it contains.
Those particles turn around a common axis while the velocity field carries and deforms the tube as a whole.
Mark off a short segment of it.
The marked particles remain the segment we follow, even as its shape changes.
When the flow draws the two ends apart, the segment lengthens along its axial direction.
The part of the velocity gradient that measures this deformation is the symmetric strain
\[
  S:=\frac12\bigl(\nabla u+(\nabla u)^{\mathsf T}\bigr),
\]
and an aligned positive direction records the lengthening directly.
If \(e\) is the axial direction and \(L(t)\) the length of the segment, then
\[
  Se=\sigma(t)e,
  \qquad
  \sigma(t)>0,
  \qquad
  \frac{\dot L(t)}{L(t)}
  =
  e\cdot Se
  =
  \sigma(t).
\]

This extension cannot occur by itself in an incompressible fluid.
The material segment keeps its volume, so the cross-section contracts while the length grows.
Let \(\mathcal V\) be its fixed material volume, set \(A(t):=\mathcal V/L(t)\), and write \(r_{\mathrm{eff}}(t):=\sqrt{A(t)/\pi}\).
Then
\[
  \nabla\cdot u=0,
  \qquad
  \frac{d}{dt}(A L)=0,
  \qquad
  \frac{\dot A}{A}=-\sigma(t),
  \qquad
  \frac{\dot r_{\mathrm{eff}}}{r_{\mathrm{eff}}}
  =
  -\frac{\sigma(t)}{2}.
\]
The segment is now longer and thinner, and both changes belong to the same motion.

The fluid inside the segment is rotating as it narrows.
Its vorticity \(\omega:=\nabla\times u\) gives the local axis and strength of that rotation.
For the viscous flow it obeys
\[
  \frac{D\omega}{Dt}
  =
  S\omega+\nu\Delta\omega,
  \qquad
  \frac{D}{Dt}
  :=
  \partial_t+u\cdot\nabla,
\]
and alignment of the vorticity with the extending direction gives
\[
  S\omega
  =
  \sigma(t)\omega,
  \qquad
  \frac12\frac{D|\omega|^2}{Dt}
  =
  \sigma(t)|\omega|^2
  +
  \nu\,\omega\cdot\Delta\omega.
\]
The stretching term amplifies the rotation already carried by the segment, while viscosity remains present in the same evolution.
On an interval where the complete right-hand side stays positive, the column becomes thinner as its rotation strengthens.

This sharpening does not require the kinetic energy of the material segment to become infinite.
Its volume is still the fixed number \(\mathcal V\), so a speed bound \(|u|\leq U_0\) on \(\Omega_{\mathrm{seg}}(t)\) gives the finite segment energy below, while the antisymmetric part of \(\nabla u\) records the growing rotation:
\[
  E_{\mathrm{seg}}(t)
  :=
  \frac12\int_{\Omega_{\mathrm{seg}}(t)}|u(x,t)|^2\,dx
  \leq
  \frac12U_0^2\mathcal V
  <
  \infty,
  \qquad
  \left|\frac12\bigl(\nabla u-(\nabla u)^{\mathsf T}\bigr)\right|_{\mathrm F}
  =
  \frac{|\omega|}{\sqrt2}
  \leq
  |\nabla u|_{\mathrm F}.
\]
The fixed volume and bounded speed keep the energy finite.
At the same time, the cross-section contracts and the rotating gradient carried through it can rise.
The change is concentrated in that narrowing cross-section, so the next comparison is made at its transverse width.

\divider{\NavierStokes{} Scaling}

Fix \(t_0\) and set \(\ell:=r_{\mathrm{eff}}(t_0)\).
The rescaled field below is a comparison solution at a finer width; it is not the original material vortex at a later time.
For \(\lambda>1\), define
\[
  u_\lambda(x,t)
  :=
  \lambda u(\lambda x,\lambda^2t),
  \qquad
  p_\lambda(x,t)
  :=
  \lambda^2p(\lambda x,\lambda^2t).
\]
At time \(\lambda^{-2}t_0\), the vortex in \(u_\lambda\) has width \(\lambda^{-1}\ell\).
Differentiating the rescaling through each term gives
\[
\begin{aligned}
  \partial_tu_\lambda(x,t)
  &=
  \lambda^3(\partial_tu)(\lambda x,\lambda^2t),
  \\
  (u_\lambda\cdot\nabla)u_\lambda(x,t)
  &=
  \lambda^3\bigl((u\cdot\nabla)u\bigr)(\lambda x,\lambda^2t),
  \\
  \nabla p_\lambda(x,t)
  &=
  \lambda^3(\nabla p)(\lambda x,\lambda^2t),
  \\
  \nu\Delta u_\lambda(x,t)
  &=
  \lambda^3\nu(\Delta u)(\lambda x,\lambda^2t).
\end{aligned}
\]
All four terms carry the same \(\lambda^3\) factor.
The finer comparison strengthens viscosity and nonlinear transport together, with pressure and local acceleration preserving the same balance.
Viscosity has gained no automatic advantage.
The equation stays balanced under concentration, but the quantities used to measure the velocity do not.

\divider{Critical Height}

Let \(\Lambda:=(-\Delta)^{1/2}\).
A change of variables shows how each homogeneous Sobolev height responds:
\[
  \norm{u_\lambda(t)}_{\dot H^s}^2
  =
  \lambda^{2s-1}
  \norm{u(\lambda^2t)}_{\dot H^s}^2.
\]
The exponent vanishes at \(s=\tfrac12\), so this height neither rises nor falls under the rescaling.
Set
\[
  H(t)
  :=
  \frac12\norm{u(t)}_{\dot H^{1/2}}^2
  =
  \frac12\norm{\Lambda^{1/2}u(t)}_2^2.
\]
With \(H_\lambda\) denoting the critical height of \(u_\lambda\), the levels on either side of it move in opposite directions:
\[
  \norm{u_\lambda(t)}_2^2
  =
  \lambda^{-1}\norm{u(\lambda^2t)}_2^2,
  \qquad
  H_\lambda(t)=H(\lambda^2t),
  \qquad
  \norm{\nabla u_\lambda(t)}_2^2
  =
  \lambda\norm{\nabla u(\lambda^2t)}_2^2.
\]
Concentration lowers kinetic energy and raises the first-derivative energy while leaving \(H\) fixed.
The critical height sits exactly where the rescaled four-term balance can be read without the measuring height moving with it.

\divider{Nonlinear Production versus Viscous Dissipation}

Return now to the evolution of the original solution.
Apply \(\Lambda^{1/2}\) to its momentum equation and pair with \(\Lambda^{1/2}u\).
The signed nonlinear production is
\[
  \mathcal P_{1/2}(t)
  :=
  -\operatorname{Re}
  \pair
    {\Lambda^{1/2}u(t)}
    {\Lambda^{1/2}\bigl((u(t)\cdot\nabla)u(t)\bigr)}_{L^2}.
\]
The pressure pairing vanishes because \(\nabla\cdot u=0\), while the Laplacian contributes \(-\nu\norm{\Lambda^{3/2}u}_2^2\).
The critical height obeys
\[
  H'(t)
  +
  \nu\norm{\Lambda^{3/2}u(t)}_2^2
  =
  \mathcal P_{1/2}(t).
\]
Thus \(H\) rises exactly on intervals where nonlinear production exceeds viscous dissipation.
Under the scale comparison, the two contributions transform as
\[
  \mathcal P_{1/2}[u_\lambda](t)
  =
  \lambda^2\mathcal P_{1/2}[u](\lambda^2t),
  \qquad
  \nu\norm{\Lambda^{3/2}u_\lambda(t)}_2^2
  =
  \lambda^2\nu\norm{\Lambda^{3/2}u(\lambda^2t)}_2^2.
\]
Production and dissipation remain matched as the width shrinks.
Their common \(\lambda^2\) factor is the inverse of the time contraction \(t\mapsto\lambda^{-2}t\): a width smaller by \(\lambda\) comes with a time shorter by \(\lambda^2\).
Viscous motion by itself makes the meaning of that clock visible.

\divider{The Heat Clock}

For the linear heat equation \(v_t=\nu\Delta v\), each Fourier mode changes according to
\[
  \widehat v(\xi,t+\delta t)
  =
  e^{-\nu|\xi|^2\delta t}\widehat v(\xi,t).
\]
A structure of width \(r\) is concentrated near frequencies \(|\xi|\simeq r^{-1}\).
Viscosity changes that structure by order one when \(\nu|\xi|^2\delta t\simeq1\), which gives
\[
  \tau_\nu(r)
  :=
  \frac{r^2}{\nu}.
\]
This is a comparison time: it measures how long diffusion needs to act substantially at width \(r\), not how long an actual \NavierStokes{} solution has been shown to take to reach that width.
At the finer width,
\[
  \tau_\nu(\lambda^{-1}r)
  =
  \lambda^{-2}\tau_\nu(r),
\]
the same square contraction carried by the full \NavierStokes{} rescaling.
Every further decrease in width brings a shorter heat clock with it.

\divider{The Zeno Cascade}

For the geometric sequence
\[
  r_n:=2^{-n}r_0,
  \qquad
  \tau_n:=\frac{r_n^2}{\nu},
\]
the clocks form the geometric series
\[
  \tau_n
  =
  \frac{r_0^2}{\nu}4^{-n},
  \qquad
  \sum_{n=0}^{\infty}\tau_n
  =
  \frac{4r_0^2}{3\nu}
  <
  \infty.
\]
Their finite sum shows only that the heat clocks do not exclude infinitely many scale transitions before a finite time.
It does not construct a singularity or show that one fluid actually makes those transitions.
For the clocks to belong to a candidate singular history, one \NavierStokes{} solution would have to pass through the widths
\[
  r_0>r_1>r_2>\cdots\longrightarrow0
\]
at times
\[
  t_0<t_1<t_2<\cdots\uparrow T_*<\infty,
  \qquad
  t_{n+1}-t_n\asymp\tau_n,
\]
with comparison constants independent of \(n\).
Its remaining time would then satisfy
\[
  T_*-t_n
  \asymp
  \sum_{k=n}^{\infty}\tau_k
  =
  \frac{4r_n^2}{3\nu}.
\]
At the \(n\)-th width, the same velocity field would have only an interval of order \(r_n^2/\nu\) in which to form the next one.
The spatial sequence has become a demand for increasingly rapid temporal response.
At a fixed spatial coordinate, the velocity must change fast enough to sharpen.
Its local acceleration has a rate of change of its own, the jerk, and every later response has another rate behind it.
Following that succession leads into the temporal Tower.
